\documentclass[11pt]{article}
\usepackage{preamble}

\newcommand{\promptplaceholder}[1]{%
  \begin{center}
  \fcolorbox{framegray!45}{frameshade}{%
    \parbox{0.84\linewidth}{\centering\small #1}%
  }
  \end{center}
}

\newcommand{\checkbox}{%
  \tikz[baseline=-0.6ex]\draw[line width=0.4pt] (0,0) rectangle (0.16,0.16);
}

\begin{document}

\packetheading{Lesson 4-4 \textperiodcentered\ Quick \textperiodcentered\ Student Packet}{Adding \& Subtracting Rational Expressions + Resistance DOK-3 (Q\#15 Prep)}

\sectionbanner{OBJECTIVES}
{\small
\renewcommand{\arraystretch}{1.25}
\noindent\begin{tabular}{|p{1.8in}|p{4.8in}|}
\hline
\textbf{Math Objective} & Add and subtract rational expressions with like and unlike denominators; identify the LCM of polynomial denominators; simplify results. \\ \hline
\textbf{Language Objective} & Explain why adding rational expressions requires common denominators, using the word `LCM' and referring to the analogy with fractions. \\ \hline
\textbf{Essential Understanding} & Rational expressions behave like rational numbers --- closed under addition and subtraction, and requiring a common denominator before you can combine them. \\ \hline
\end{tabular}
}

\sectionbanner{TOPIC GOALS / ESSENTIAL QUESTION / MATERIALS}
{\small
\renewcommand{\arraystretch}{1.25}
\noindent\begin{tabular}{|p{1.8in}|p{4.8in}|}
\hline
\textbf{Topic Goals} & Fluency adding/subtracting rational expressions; understanding closure under these operations (Topic 4 assessment Q\#15 prep). \\ \hline
\textbf{Essential Question} & When you add two rational expressions with different denominators, why must you find the LCM first? \\ \hline
\textbf{Materials} & Student packet, pencil, calculator. DOK-3 Practice \#32 (electrical resistance) is the capstone and rehearses the structural reasoning tested on Topic 4 LEHS Q\#15. \\ \hline
\end{tabular}
}

\sectionbanner{LAUNCH --- Adding \& Subtracting Rational Expressions (teacher models Ex 3 + Ex 4)}

\begin{callout}{\IconBook\ RULES --- ADD/SUBTRACT RATIONAL EXPRESSIONS (keep printed on your page)}{calloutgreen}
\begin{enumerate}[leftmargin=2.1em,itemsep=1pt,topsep=3pt]
  \item LIKE denominators: combine numerators, keep the denominator, simplify.
  \item UNLIKE denominators:
  \begin{enumerate}[label=(\alph*),leftmargin=1.8em,itemsep=1pt,topsep=2pt]
    \item factor each denominator,
    \item find the LCM,
    \item rewrite each fraction with the LCM as denominator,
    \item combine numerators,
    \item simplify.
  \end{enumerate}
  \item The LCM of two polynomial denominators includes every DISTINCT factor, taken to its highest power across both.
  \item Always check whether the final expression can be simplified by cancelling a common factor.
\end{enumerate}
\end{callout}

\begin{callout}{\IconWarn\ COMPRESSED LESSON --- TWO EXAMPLES IN LAUNCH}{calloutred}
This is a single-period quick treatment. Teacher models BOTH Example 3 (add unlike denominators) AND Example 4 (subtract unlike denominators) during Launch.\\
Watch carefully on Example 4 --- when you subtract, you must distribute the minus sign across every term of the second numerator. This is the \#1 error source.
\end{callout}

\sectionbanner{EXPLORE --- Think-Pair-Share (32 min)}
\textit{Work with a partner. Move in order. Practice \#32 (Electrical Resistance) is the big one --- plan \(\sim\)10 min for it.}\\
\textit{45-min Wed F variant: skip \#16 and \#26.}

\textbf{Bank-item labels (in order):}
\begin{enumerate}[leftmargin=1.6em,itemsep=2pt,topsep=2pt]
  \item Try It 3 --- Add rational expressions (unlike denominators)
  \item Try It 4 --- Subtract rational expressions (unlike denominators)
  \item Practice \#12 --- Like denominators
  \item Practice \#16 --- Unlike denominators, factor first
  \item Practice \#26 --- Subtract with sign-distribution trap
  \item Practice \#32 \IconStar\ DOK-3 ELECTRICAL RESISTANCE --- Topic 4 LEHS Q\#15 prep
\end{enumerate}

\bankitem{Try It 3 --- Add rational expressions (unlike denominators)}{%
Find the sum.

\begin{enumerate}[label=\textbf{\alph*.},leftmargin=1.8em,itemsep=3pt,topsep=4pt]
  \item \(\displaystyle \frac{x+6}{x^2-4}+\frac{2}{x^2-5x+6}\)
  \item \(\displaystyle \frac{2x}{3x+4}+\frac{4x^2-11x-12}{6x^2+5x-4}\)
\end{enumerate}
}{}

\bankitem{Try It 4 --- Subtract rational expressions (unlike denominators)}{%
Simplify.

\begin{enumerate}[label=\textbf{\alph*.},leftmargin=1.8em,itemsep=3pt,topsep=4pt]
  \item \(\displaystyle \frac{1}{3x}+\frac{1}{6x}-\frac{1}{x^2}\)
  \item \(\displaystyle \frac{3x-5}{x^2-25}-\frac{2}{x+5}\)
\end{enumerate}
}{}

\bankitem{Practice \#12 --- Like denominators}{%
\textbf{Generalize} Explain how addition and subtraction of rational expressions is similar to and different from addition and subtraction of rational numbers.
}{}

\bankitem{Practice \#16 --- Unlike denominators, factor first}{%
\textbf{Look for Relationships} Find the LCM of \(3x^2 + 7x + 2\) and \(9x + 3\).
}{}

\bankitem{Practice \#26 --- Subtract with sign-distribution trap}{%
On Saturday morning, Ahmed decided to take a bike ride from one end of the 15-mile bike trail to the other end of the bike trail and back. His average speed the first half of the ride was 2 mph faster than his speed the second half. Find an expression for Ahmed's total travel time. If his average speed for the first half of the ride was 12 mph, how long was Ahmed's bike ride?

\promptplaceholder{[IMAGE: Bike trail with cyclist, sign "Trail length 15 miles", one label "Speed x+2" and another label "Speed x".]}
}{}

\bankitem{Practice \#32 \IconStar\ DOK-3 ELECTRICAL RESISTANCE --- Topic 4 LEHS Q\#15 prep}{%
\textbf{Model With Mathematics} The resistance of electrical circuits A and B are shown. Find the ratio of the resistance of Circuit A to Circuit B.

\promptplaceholder{[GRAPH / TIKZ figure]}
}{}

\sectionbanner{SHARE / SUMMARY (5 min)}

\begin{sentenceframebox}
\textbf{Sentence frames}

Pair share (Resistance): ``For the parallel circuit, \(1/R_{\text{total}} = 1/R_1 + 1/R_2\). I found a common denominator of \_\_\_, added the reciprocals to get \(1/R_{\text{total}} =\) \_\_\_, then flipped to get \(R_{\text{total}} =\) \_\_\_. The ratio \(R_A/R_B\) simplified to \_\_\_ because \_\_\_.'' 

Whole class: one pair reads their resistance solution. Class signals thumbs up/sideways.
\end{sentenceframebox}

\begin{callout}{\IconSearch\ BRIDGE TO TOPIC 4 LEHS Q\#15}{calloutpeach}
The resistance move --- find LCM, combine reciprocals into a single rational expression --- is the structural reasoning tested on Topic 4 LEHS Q\#15 (operations/closure). You have now rehearsed it.
\end{callout}

\sectionbanner{EXIT}

\begin{summaryexitbox}{EXIT TICKET --- Summary Recap}
To add rational expressions with unlike denominators, I must first find the \_\_\_\_\_ of the denominators.

Practice \#32 (resistance) showed me that reciprocal-sum problems (like \(1/R_1 + 1/R_2\)) use the same algebraic steps as adding rational expressions because \_\_\_\_\_.

One biggest thing I learned today: \_\_\_\_\_\_\_\_\_\_\_\_\_\_\_\_\_\_\_\_\_\_\_\_
\end{summaryexitbox}

\clearpage

\sectionbanner{OPTIONAL REINFORCEMENT (not collected)}
\textit{If you finish Explore early. \#34 is a compound-fraction equivalence multi-select that rehearses closure under addition and subtraction.}

\bankitem{Optional \#1  (Savvas Practice \#34 --- compound-fraction multi-select)}{%
Which of the following compound fractions simplifies to \(\frac{x + 1}{x - 3}\)? Select all that apply.

\begin{itemize}[leftmargin=1.8em,itemsep=4pt,topsep=4pt]
  \item[\checkbox] \(\displaystyle \frac{\frac{x^2 + 5x + 4}{x^2 + 2x - 8}}{\frac{x^2 - 4x + 3}{x^2 - 3x + 2}}\)
  \item[\checkbox] \(\displaystyle \frac{\frac{x^2 - 1}{x^2 - 4}}{\frac{x^2 + x - 7}{x^2 + 5x + 6}}\)
  \item[\checkbox] \(\displaystyle \frac{\frac{x^2 + 3x - 10}{x^2 - 16}}{\frac{x^2 - 4x - 5}{x^2 - 1}}\)
  \item[\checkbox] \(\displaystyle \frac{\frac{x^2 + 3x - 10}{x^2 - 5x + 6}}{\frac{x^2 - 25}{x^2 - 4x - 5}}\)
\end{itemize}
}{}

\end{document}
